
\SGARefTarget{sga3:iv:section:6}\section*{6. Topologies in the Category of Schemes}
\addcontentsline{toc}{section}{6. Topologies in the Category of Schemes}

\SGARefTarget{sga3:iv:subsection:6:1}\subsection*{6.1. The Zariski Topology}
\addcontentsline{toc}{subsection}{\SGARefLink{sga3:iv:subsection:6:1}{6.1}. The Zariski Topology}

This is the topology generated by the following pretopology: a family of
morphisms \(\{S_i\to S\}\) is covering if each morphism is an open
immersion and the union of the images of the \(S_i\) is all of \(S\). It is
denoted by \((\mathrm{Zar})\).

\medskip
\noindent\SGARefTarget{sga3:iv:definition:6:1:1}\textbf{Definition 6.1.1.}
\textit{A sheaf for the Zariski topology is also called a \emph{functor of
local nature}: it is a contravariant functor from \(\mathbf{Sch}\) to
\(\mathbf{Set}\) such that, for every scheme \(S\) and every covering of
\(S\) by open subschemes \(S_i\), the diagram}
\[
  F(S)\longrightarrow\prod_iF(S_i)
  \rightrightarrows\prod_{i,j}F(S_i\cap S_j)
\]
\textit{is exact.}

In particular, a functor of local nature carries direct sums to products.
Since every representable functor is a sheaf, the Zariski topology is
coarser than the canonical topology.

As a matter of terminology, whenever we say "local" or "locally" without
further qualification, the reference is to the Zariski topology and hence
to the usual meaning of these words.

\SGARefTarget{sga3:iv:subsection:6:2}\subsection*{6.2. A Procedure for Constructing Topologies}
\addcontentsline{toc}{subsection}{\SGARefLink{sga3:iv:subsection:6:2}{6.2}. A Procedure for Constructing Topologies}

\medskip
\noindent\SGARefTarget{sga3:iv:proposition:6:2:1}\textbf{Proposition 6.2.1.}
\textit{Let \(\mathscr C\) be a category, let \(\mathscr C'\) be a full
subcategory, and let \(P\) be a set of families of morphisms in
\(\mathscr C\) with common target, stable under base change and composition
(that is, satisfying axioms \SGARefLink{sga3:iv:definition:4:2:5:item:p:1:01}{(P~1)} and \SGARefLink{sga3:iv:definition:4:2:5:item:p:2:01}{(P~2)} of \SGARefLink{sga3:iv:definition:4:2:5}{4.2.5}). Let \(P'\) be a set
of families of morphisms in \(\mathscr C'\) that contains every singleton
family consisting of an identity isomorphism. Equip \(\mathscr C\) with
the topology generated by \(P\) and \(P'\) (cf. \SGARefLink{sga3:iv:definition:4:2:5:0}{4.2.5.0}), and suppose that
the following three conditions hold:}

\begin{enumerate}
  \item[(a)]\SGARefTarget{sga3:iv:proposition:6:2:1:item:a:01} \textit{If \(\{S_i\to S\}\in P'\), so that
    \(S_i,S\in\operatorname{Ob}\mathscr C'\), and \(T\to S\) is a
    morphism in \(\mathscr C'\), then the fiber products
    \(S_i\times_ST\), formed in \(\mathscr C\), exist and the family
    \(\{S_i\times_ST\to T\}\) belongs to \(P'\); in particular,
    \(S_i\times_ST\in\operatorname{Ob}\mathscr C'\).}

    \textit{\SGARefLink{sga3:iv:proposition:6:2:1:item:a:01}{This condition} implies that \(P'\) is stable under base change
    in \(\mathscr C'\), but is not equivalent to that assertion: it also
    requires the inclusion functor \(\mathscr C'\hookrightarrow\mathscr C\)
    to commute with certain fiber products.}

  \item[(b)]\SGARefTarget{sga3:iv:proposition:6:2:1:item:b:01} \textit{For every \(S\in\operatorname{Ob}\mathscr C\), there
    exists a family \(\{S_i\to S\}\in P\) with
    \(S_i\in\operatorname{Ob}\mathscr C'\) for every \(i\).}

  \item[(c)]\SGARefTarget{sga3:iv:proposition:6:2:1:item:c:01} \textit{In a situation}
    \[
    \begin{array}{ccccc}
      & & S_{ijk} & &\\
      & & \scriptstyle(P')\,\Big\downarrow & &\\[-0.25em]
      S_i & \xleftarrow{\ (P)\ } & S_{ij} & &\\
      \scriptstyle(P')\,\Big\downarrow & & & &\\[-0.25em]
      S & & & &
    \end{array}
    \tag{6.2.1.c}\label{sga3:iv:proposition:6:2:1:eq:6:2:1:c:01}
    \]
    \textit{where \(S,S_i,S_{ij},S_{ijk}\) are objects of
    \(\mathscr C'\), where \(\{S_i\to S\}\in P'\),
    \(\{S_{ij}\to S_i\}\in P\) for each \(i\), and
    \(\{S_{ijk}\to S_{ij}\}\in P'\) for each pair \(ij\), there exists a
    family \(\{T_n\to S\}\in P'\), and for each \(n\) a multi-index
    \((ijk)\), fitting into a commutative diagram}
    \[
    \begin{array}{ccccc}
      S_{ijk} & \longleftarrow & T_n\\
      & \searrow & \Big\downarrow\\[-0.25em]
      & & S .
    \end{array}
    \tag{6.2.1.d}\label{sga3:iv:proposition:6:2:1:eq:6:2:1:d:01}
    \]
\end{enumerate}

\textit{Under these hypotheses, a sieve \(R\) of an object
\(S\in\operatorname{Ob}\mathscr C\) is covering if and only if there is a
composite family}
\[
\begin{array}{ccccc}
  R & \dashleftarrow & S_{ij}\\
  \Big\downarrow & & \scriptstyle(P')\,\Big\downarrow\\[-0.25em]
  S & \xleftarrow{\ (P)\ } & S_i
\end{array}
\tag{6.2.1.e}\label{sga3:iv:proposition:6:2:1:eq:6:2:1:e:01}
\]
\textit{where \(S_i,S_{ij}\in\operatorname{Ob}\mathscr C'\),
\(\{S_i\to S\}\in P\), and \(\{S_{ij}\to S_i\}\in P'\) for every \(i\),
such that the resulting morphisms \(S_{ij}\to S\) factor through \(R\).
Equivalently, the sieve generated by this composite family is contained in
\(R\).}

\smallskip
\noindent\textit{Proof.}
Families belonging to \(P\) and \(P'\) are covering, and therefore so is a
composite of such families, by \SGARefLink{sga3:iv:definition:4:2:3:item:c:2:01}{(C~2)}. A sieve of the displayed form is
thus covering because it contains a covering sieve.

Conversely, it is enough to show that the sieves of this form constitute a
topology, that is, to verify axioms \SGARefLink{sga3:iv:definition:4:2:1:item:t:1:01}{(T~1)}--\SGARefLink{sga3:iv:definition:4:2:1:item:t:4:01}{(T~4)} of \SGARefLink{sga3:iv:definition:4:2:1}{4.2.1}.

\smallskip
\noindent\textit{Axiom \SGARefLink{sga3:iv:definition:4:2:1:item:t:4:01}{(T~4)}.}
Let \(S\in\operatorname{Ob}\mathscr C\). By \SGARefLink{sga3:iv:proposition:6:2:1:item:b:01}{(b)}, there exists a family
\(\{S_i\to S\}\in P\) with
\(S_i\in\operatorname{Ob}\mathscr C'\). The singleton families
\(\{S_i\xrightarrow{\operatorname{id}_{S_i}}S_i\}\) belong to \(P'\) by
hypothesis. Hence the maximal sieve \(S\) of \(S\) has the required form:
\[
\begin{array}{ccc}
  S_i & \longrightarrow & S\\
  \scriptstyle(P')\,\Big\downarrow\rlap{\(\scriptstyle\operatorname{id}\)}
      & & \Big\downarrow\rlap{\(\scriptstyle\operatorname{id}_S\)}\\[-0.25em]
  S_i & \xrightarrow{\ (P)\ } & S .
\end{array}
\tag{6.2.1.T4}\label{sga3:iv:proposition:6:2:1:eq:6:2:1:t4:01}
\]

\smallskip
\noindent\textit{Axiom \SGARefLink{sga3:iv:definition:4:2:1:item:t:3:01}{(T~3)}.}
This is evident.

\smallskip
\noindent\textit{Axiom \SGARefLink{sga3:iv:definition:4:2:1:item:t:2:01}{(T~2)}.}
Let \(R\) be a sieve of \(S\) of the required form, and let \(C\) be a
sieve\footnote{Editor's note: the original has been corrected here.} of
\(S\) such that, for every \(T\in\operatorname{Ob}\mathscr C\) and every
morphism \(T\to S\) factoring through \(R\), the sieve \(C\times_ST\) of
\(T\) has the required form. Since \(S_{ij}\to S\) factors through \(R\),
the sieve
\[
  C_{ij}=C\times_SS_{ij}
\]
of \(S_{ij}\) has the required form; it is represented schematically by
\[
\begin{array}{ccccc}
    & & S_{ij} & \longleftarrow & C_{ij}\\
    & \swarrow & \scriptstyle(P')\,\Big\downarrow & &\\[-0.25em]
    & & S_i & & \Bigg\downarrow\\[-0.25em]
    & & \scriptstyle(P)\,\Big\downarrow & &\\[-0.25em]
  R & \longrightarrow & S & \longleftarrow & C .
\end{array}
\tag{6.2.1.T2a}\label{sga3:iv:proposition:6:2:1:eq:6:2:1:t2a:01}
\]
Consequently, for every pair \(ij\), there is a diagram of the form
\[
\begin{array}{ccccc}
  S_{ijkl} & & & &\\
  \scriptstyle(P')\,\Big\downarrow & \searrow & & &\\[-0.25em]
  S_{ijk} & & & &\\
  \scriptstyle(P)\,\Big\downarrow & & & &\\[-0.25em]
  S_{ij} & \longleftarrow & C_{ij}. & &
\end{array}
\tag{6.2.1.T2b}\label{sga3:iv:proposition:6:2:1:eq:6:2:1:t2b:01}
\]
We have therefore produced a composite family
\[
  S_{ijkl}\xrightarrow{\ (P')\ }S_{ijk}
  \xrightarrow{\ (P)\ }S_{ij}
  \xrightarrow{\ (P')\ }S_i
  \xrightarrow{\ (P)\ }S
\tag{6.2.1.T2c}\label{sga3:iv:proposition:6:2:1:eq:6:2:1:t2c:01}
\]
belonging to \(P\circ P'\circ P\circ P'\), factoring through \(C\), and
having every object other than \(S\) in \(\mathscr C'\).

Apply condition \SGARefLink{sga3:iv:proposition:6:2:1:item:c:01}{(c)} to each family \(\{S_{ijkl}\to S_i\}\). For every \(i\)
we obtain a family \(\{T_{in}\to S_i\}\in P'\) such that \(T_{in}\to S\)
factors through one of the \(S_{ijkl}\), and hence through \(C\):
\[
\begin{array}{ccccc}
  T_{in} & \longrightarrow & S_{ijkl}\\
  \scriptstyle(P')\,\Big\downarrow & & \Big\downarrow\\[-0.25em]
  S_i & & C\\
  \scriptstyle(P)\,\Big\downarrow & \swarrow &\\[-0.25em]
  S & &
\end{array}
\tag{6.2.1.T2d}\label{sga3:iv:proposition:6:2:1:eq:6:2:1:t2d:01}
\]
Thus the sieve \(C\) of \(S\) has the required form, completing the
verification of \SGARefLink{sga3:iv:definition:4:2:1:item:t:2:01}{(T~2)}.

\smallskip
\noindent\textit{Axiom \SGARefLink{sga3:iv:definition:4:2:1:item:t:1:01}{(T~1)}.}
Let \(R\) be a sieve of \(S\) of the prescribed form and let \(T\to S\) be
a morphism in \(\mathscr C\). We show that the sieve \(R\times_ST\) of
\(T\) again has the prescribed form. Consider the commutative configuration
\[
\begin{array}{ccccccc}
  & & S_{ij} & \multicolumn{3}{c}{\xleftarrow{\hspace{6em}}} & U_{ikj}\\
  & \swarrow & \scriptstyle(P')\,\Big\downarrow & & & &
    \scriptstyle(P')\,\Big\downarrow\\[-0.25em]
  & & S_i & \longleftarrow & T_i & \xleftarrow{\ (P)\ } & U_{ik}\\
  & & \scriptstyle(P)\,\Big\downarrow & &
    \scriptstyle(P)\,\Big\downarrow & \swarrow_{\scriptstyle(P)} &\\[-0.25em]
  R & \longrightarrow & S & \longleftarrow & T & & .
\end{array}
\tag{6.2.1.T1a}\label{sga3:iv:proposition:6:2:1:eq:6:2:1:t1a:01}
\]
Put \(T_i=S_i\times_ST\). The family \(\{T_i\to T\}\) belongs to \(P\),
by \SGARefLink{sga3:iv:definition:4:2:5:item:p:1:01}{(P~1)}. By \SGARefLink{sga3:iv:proposition:6:2:1:item:b:01}{(b)}, choose families \(\{U_{ik}\to T_i\}\in P\) with
\(U_{ik}\in\operatorname{Ob}\mathscr C'\). By \SGARefLink{sga3:iv:definition:4:2:5:item:p:2:01}{(P~2)},
\(\{U_{ik}\to T\}\in P\). Condition \SGARefLink{sga3:iv:proposition:6:2:1:item:a:01}{(a)} shows that
\[
  U_{ik}\times_{S_i}S_{ij}=U_{ikj}
\]
is an object of \(\mathscr C'\) and that, for each \(ik\),
\(\{U_{ikj}\to U_{ik}\}\in P'\). The commutative diagram
\[
\begin{array}{ccccc}
  R & \longleftarrow & U_{ikj}\\
  \Big\downarrow & & \scriptstyle(P')\,\Big\downarrow\\[-0.25em]
  S & \longleftarrow & T & \xleftarrow{\ (P)\ } U_{ik}
\end{array}
\tag{6.2.1.T1b}\label{sga3:iv:proposition:6:2:1:eq:6:2:1:t1b:01}
\]
shows that the morphisms \(U_{ikj}\to T\) factor through the sieve
\(T\times_SR\) of \(T\). This sieve therefore has the required form, and
the proof is complete.

\medskip
\noindent\SGARefTarget{sga3:iv:corollary:6:2:2}\textbf{Corollary 6.2.2.}
\textit{Let \(S\in\operatorname{Ob}\mathscr C'\), and let \(R\) be a
sieve of \(S\). Then \(R\) is covering if and only if there exists a family
\(\{T_i\to S\}\in P'\) that factors through \(R\).}

Such a sieve is plainly covering. Conversely, apply condition \SGARefLink{sga3:iv:proposition:6:2:1:item:c:01}{(c)} with the
family \(\{S_i\to S\}\) reduced to the identity isomorphism of \(S\); \SGARefLink{sga3:iv:proposition:6:2:1}{the
proposition} then shows that every covering sieve has the stated form.

\medskip
\noindent\SGARefTarget{sga3:iv:corollary:6:2:3}\textbf{Corollary 6.2.3.}
\textit{A presheaf \(F\) on \(\mathscr C\) is separated (respectively, is
a sheaf) if and only if the map}
\[
  F(S)\longrightarrow\prod_iF(S_i)
\]
\textit{is injective (respectively, the diagram}
\[
  F(S)\longrightarrow\prod_iF(S_i)
  \rightrightarrows\prod_{i,j}F(S_i\times_SS_j)
\]
\textit{is exact) in each of the following two cases:}
\begin{enumerate}
  \item[(i)]\SGARefTarget{sga3:iv:corollary:6:2:3:item:i:01} \textit{\(\{S_i\to S\}\in P\);}
  \item[(ii)]\SGARefTarget{sga3:iv:corollary:6:2:3:item:ii:01} \textit{\(S,S_i\in\operatorname{Ob}\mathscr C'\) and
    \(\{S_i\to S\}\in P'\).}
\end{enumerate}

The conditions are necessary because the families in question are
covering. Let \(C\) be the sieve of \(S\) that is the image of a composite
family
\[
  \{S_{ij}\xrightarrow{\ (P')\ }S_i\xrightarrow{\ (P)\ }S\}.
\]
A straightforward diagram chase shows that the conditions of \SGARefLink{sga3:iv:corollary:6:2:3}{the
corollary} imply that
\[
  \operatorname{Hom}(S,F)
  \xrightarrow{\ g\ }\operatorname{Hom}(C,F)
\]
is injective, respectively bijective. Every covering sieve \(R\) of \(S\)
contains a sieve \(C\) of this form, and there is a commutative diagram
\[
\begin{array}{ccccc}
  \operatorname{Hom}(S,F) & \xrightarrow{\ f\ } &
    \operatorname{Hom}(R,F)\\
  & \searrow_{\scriptstyle g} & \Big\downarrow\rlap{\(\scriptstyle h\)}\\[-0.25em]
  & & \operatorname{Hom}(C,F).
\end{array}
\tag{6.2.3.1}\label{sga3:iv:corollary:6:2:3:eq:6:2:3:1:01}
\]
Since \(g\) is injective, so is \(f\), and hence \(F\) is separated.
Moreover, \(R\) refines \(C\), so \(h\) is also injective. If \(g\) is
bijective, then \(f\) is bijective as well, and \(F\) is a sheaf.

\medskip
\noindent\SGARefTarget{sga3:iv:remark:6:2:4}\textbf{Remark 6.2.4.}
\textit{\SGARefLink{sga3:iv:corollary:6:2:3}{The preceding corollary} does not follow from \SGARefLink{sga3:iv:corollary:4:3:5}{4.3.5}, since \(P'\) is
not stable under base extension.}

\medskip
\noindent\SGARefTarget{sga3:iv:remark:6:2:5}\textbf{Remark 6.2.5.}
\textit{Condition \SGARefLink{sga3:iv:proposition:6:2:1:item:c:01}{(c)} holds in particular when:}
\begin{enumerate}
  \item[(i)]\SGARefTarget{sga3:iv:remark:6:2:5:item:i:01} \textit{\(P'\) is stable under composition;}
  \item[(ii)]\SGARefTarget{sga3:iv:remark:6:2:5:item:ii:01} \textit{whenever \(\{S_j\to S\}\) is a family of morphisms in
    \(\mathscr C'\) belonging to \(P\), it contains a subfamily belonging
    to \(P'\).}
\end{enumerate}

\SGARefTarget{sga3:iv:subsection:6:3}\subsection*{6.3. Application to the Category of Schemes}
\addcontentsline{toc}{subsection}{\SGARefLink{sga3:iv:subsection:6:3}{6.3}. Application to the Category of Schemes}

Take \(\mathscr C\) to be the category of schemes, \(\mathscr C'\) the full
subcategory of affine schemes, and \(P\) the set of surjective families of
open immersions. We shall consider the following sets \(P'_i\):
\begin{align*}
 P'_1 &: \text{finite surjective families of flat morphisms};\\
 P'_2 &: \text{finite surjective families of flat, finitely presented,
          quasi-finite morphisms};\\
 P'_3 &: \text{finite surjective families of \'{e}tale morphisms};\\
 P'_4 &: \text{finite surjective families of finite \'{e}tale morphisms}.
\end{align*}
The description of \(P'_2\) is accompanied by the following
clarification.\footnote{Editor's note: let \(P''_2\) be the set of finite
surjective families of morphisms in \(\mathscr C'\) that are flat and of
finite presentation. By Proposition~\SGARefLink{sga3:iv:proposition:6:3:1}{6.3.1} below, the topology \(T_2\)
generated by \(P\) and \(P'_2\) coincides with the topology generated by
\(P\) and \(P''_2\). This follows from the results of
EGA~IV\textsubscript{4}, \S17.16 on quasi-sections; see the proof of \SGARefLink{sga3:iv:proposition:6:3:1}{6.3.1}.
We recall the following special case of EGA~IV\textsubscript{4}, 17.16.2.
Let \(S\) be affine and let \(f:X\to S\) be a surjective flat morphism
locally of finite presentation. Then there exist a faithfully flat,
finitely presented, quasi-finite morphism \(S'\to S\), with \(S'\) affine,
and an \(S\)-morphism \(S'\to X\).}

For each \(P'_i\), except \(P'_4\), the conditions of Proposition~\SGARefLink{sga3:iv:proposition:6:2:1}{6.2.1}
hold. Condition \SGARefLink{sga3:iv:proposition:6:2:1:item:c:01}{(c)} follows from \SGARefLink{sga3:iv:remark:6:2:5}{6.2.5}: since an affine scheme is
quasi-compact, every family of morphisms in \(\mathscr C'\) belonging to
\(P\) contains a finite subfamily that still belongs to \(P\), and hence
belongs to \(P'_i\) for \(i=1,2,3\). The topology \(T_i\) generated by
\(P\) and \(P'_i\) is denoted and named as follows:\footnotemark[\value{footnote}]
\[
\begin{array}{ccl}
 T_1&=(\mathrm{fpqc})&
   \text{the faithfully flat quasi-compact topology},\\
 T_2&=(\mathrm{fppf})&
   \text{the faithfully flat topology (locally of finite presentation)},\\
 T_3&=(\mathrm{\acute et})&
   \text{the \'{e}tale topology},\\
 T_4&=(\mathrm{\acute etf})&
   \text{the finite \'{e}tale topology}.
\end{array}
\]
Since \(P'_1\supset P'_2\supset P'_3\supset P'_4\), one has
\[
  (\mathrm{fpqc})\geq(\mathrm{fppf})\geq(\mathrm{\acute et})
  \geq(\mathrm{\acute etf})\geq(\mathrm{Zar}).
\]

\medskip
\noindent\SGARefTarget{sga3:iv:proposition:6:3:1}\textbf{Proposition 6.3.1.}
\begin{enumerate}
  \item[(i)]\SGARefTarget{sga3:iv:proposition:6:3:1:item:i:01} \textit{A sieve \(R\) of \(S\) is covering for \(T_i\),
    \(1\leq i\leq3\), if and only if there exist a covering
    \((S_p)\) of \(S\) by affine open subschemes and, for every \(p\), a
    family \(\{S_{pq}\to S_p\}\in P'_i\), with all \(S_{pq}\) affine,
    such that every morphism \(S_{pq}\to S\) factors through \(R\).}
    \footnote{Editor's note: every family
    \(\{S_{pq}\to S_p\}\in P'_i\) is finite. Hence
    \(S'_p=\coprod_qS_{pq}\) is affine, and the family may be replaced by
    the morphism \(S'_p\to S_p\), which still belongs to \(P'_i\).}

  \item[(ii)]\SGARefTarget{sga3:iv:proposition:6:3:1:item:ii:01} \textit{A presheaf \(F\) on \(\mathbf{Sch}\) is a sheaf for
    \((\mathrm{fpqc})\) (respectively, \((\mathrm{fppf})\),
    \((\mathrm{\acute et})\), or \((\mathrm{\acute etf})\)) if and only
    if:}
    \begin{enumerate}
      \item[(a)]\SGARefTarget{sga3:iv:proposition:6:3:1:item:a:01} \textit{\(F\) is a sheaf for \((\mathrm{Zar})\), that is,
        a functor of local nature;}
      \item[(b)]\SGARefTarget{sga3:iv:proposition:6:3:1:item:b:01} \textit{for every faithfully flat (respectively,
        faithfully flat, finitely presented and quasi-finite; surjective
        \'{e}tale; or surjective finite \'{e}tale) morphism \(T\to S\),
        with \(T\) and \(S\) affine, the diagram}
        \[
          F(S)\longrightarrow F(T)\rightrightarrows F(T\times_ST)
        \]
        \textit{is exact.}
    \end{enumerate}

  \item[(iii)]\SGARefTarget{sga3:iv:proposition:6:3:1:item:iii:01} \textit{The topologies \(T_i\), for \(i=1,2,3,4\), are
    coarser than the canonical topology.}

  \item[(iv)]\SGARefTarget{sga3:iv:proposition:6:3:1:item:iv:01} \textit{Every surjective family of flat open morphisms
    (respectively, flat morphisms locally of finite presentation;
    \'{e}tale morphisms; or finite \'{e}tale morphisms) is covering for
    \((\mathrm{fpqc})\) (respectively, \((\mathrm{fppf})\),
    \((\mathrm{\acute et})\), or \((\mathrm{\acute etf})\)).}

  \item[(v)]\SGARefTarget{sga3:iv:proposition:6:3:1:item:v:01} \textit{Every finite surjective family of flat quasi-compact
    morphisms is covering for \((\mathrm{fpqc})\). In particular, every
    faithfully flat quasi-compact morphism is covering for
    \((\mathrm{fpqc})\).}
\end{enumerate}

\smallskip
\noindent\textit{Proof.}
Part \SGARefLink{sga3:iv:proposition:6:3:1:item:i:01}{(i)} follows from \SGARefLink{sga3:iv:proposition:6:2:1}{6.2.1} and part \SGARefLink{sga3:iv:proposition:6:3:1:item:ii:01}{(ii)} from \SGARefLink{sga3:iv:corollary:6:2:3}{6.2.3}, taking into account
that a sheaf for the Zariski topology carries direct sums to products.
Every representable functor is a sheaf for \((\mathrm{Zar})\) and satisfies
condition \SGARefLink{sga3:iv:proposition:6:3:1:item:ii:01}{(ii)}\SGARefLink{sga3:iv:proposition:6:3:1:item:b:01}{(b)}, by \SGARefLink{sga3:iv:bib:sga:1}{SGA~1}, VIII, 5.3. Thus \(T_1\) is coarser than the
canonical topology, which proves \SGARefLink{sga3:iv:proposition:6:3:1:item:iii:01}{(iii)}.

We prove \SGARefLink{sga3:iv:proposition:6:3:1:item:iv:01}{(iv)}. Let \(\{S_i\to S\}\) be a family as in the statement. By
covering \(S\) with affine open subschemes, we reduce immediately to the
case in which \(S\) is affine.\footnote{Editor's note: the argument that
follows has been simplified by using the fact that \(S\) is now assumed
affine.}

First suppose that the morphisms \(S_i\to S\) are flat and open
(respectively, \'{e}tale). Cover each \(S_i\) by affine open subschemes
\(S_{ij}\). Since the morphisms are open, the images \(T_{ij}\) of the
\(S_{ij}\) form an open covering of \(S\). The affine scheme \(S\) is
quasi-compact, so finitely many \(T_{ij}\), indexed by a finite set \(E\),
cover it. Then
\[
  S'=\coprod_{(i,j)\in E}S_{ij}
\]
is affine, and \(S'\to S\) belongs to \(P'_1\), respectively \(P'_3\), and
is therefore covering. Since it factors through the given family, that
family is covering.

In the \((\mathrm{\acute etf})\) case, every \(S_i\) is finite over \(S\)
and hence affine. In the preceding argument one may take the singleton
covering \(\{S_i\}\) of \(S_i\), which gives \(S'\to S\in P'_4\).

Now suppose that the morphisms \(f_i:S_i\to S\) are flat and locally of
finite presentation. For every \(s\in S\), the proof of
EGA~IV\textsubscript{4}, 17.16.2 supplies an affine subscheme \(X(s)\) of
some \(S_i\), with \(s\in f_i(X(s))\), such that the restriction
\[
  g_i:X(s)\longrightarrow S
\]
is flat, finitely presented, and quasi-finite. Its image \(g_i(X(s))\) is
an open neighborhood \(U(s)\) of \(s\), by EGA~IV\textsubscript{2}, 2.4.6.
Since \(S\) is affine, finitely many such opens
\(U(s_j)\), \(j=1,\ldots,n\), cover \(S\). Consequently,
\[
  X'=\coprod_{j=1}^{n}X(s_j)
\]
is affine and \(X'\to S\) is surjective, flat, finitely presented, and
quasi-finite, so it belongs to \(P'_2\).\footnote{Editor's note: this shows
that if \(P''_2\) denotes the set of finite surjective families of flat,
finitely presented morphisms in \(\mathscr C'\), then the topology
generated by \(P\) and \(P''_2\) equals \(T_2\). Moreover, in the notation
used at the beginning of the proof of \SGARefLink{sga3:iv:proposition:6:3:1:item:iv:01}{(iv)}, if each \(S_i\) is covered by
affine opens of finite presentation over \(S\), then the resulting
morphism \(S'\to S\) belongs to \(P''_2\).}
This completes the proof of \SGARefLink{sga3:iv:proposition:6:3:1:item:iv:01}{(iv)}.

We prove \SGARefLink{sga3:iv:proposition:6:3:1:item:v:01}{(v)}. Let \(\{S_i\to S\}\) be a finite surjective family of flat
quasi-compact morphisms, and cover \(S\) by affine open subschemes \(T_j\).
The schemes
\[
  S_{ij}=T_j\times_SS_i
\]
are quasi-compact, and therefore admit finite affine open coverings
\(T_{ijk}\). Each \(T_{ijk}\to T_j\) is flat, and the family
\(\{T_{ijk}\to T_j\}\) is finite and surjective, hence covering for \(T_1\).
By composition, \(\{T_{ijk}\to S\}\) is covering as well. It factors
through the given family, which is therefore covering:
\[
\begin{array}{ccccc}
  S_i & \longleftarrow & S_{ij} & \longleftarrow & T_{ijk}\\
  \Big\downarrow & & \Big\downarrow & \swarrow &\\[-0.25em]
  S & \longleftarrow & T_j & &
\end{array}
\tag{6.3.1.v}\label{sga3:iv:proposition:6:3:1:eq:6:3:1:v:01}
\]

\medskip
\noindent\SGARefTarget{sga3:iv:corollary:6:3:2}\textbf{Corollary 6.3.2.}
\textit{Let \((M_i)\) be the following families of morphisms:}
\begin{align*}
 (M_1)&:\ \text{faithfully flat quasi-compact morphisms};\\
 (M_2)&:\ \text{faithfully flat morphisms locally of finite presentation};\\
 (M_3)&:\ \text{surjective \'{e}tale morphisms};\\
 (M_4)&:\ \text{surjective finite \'{e}tale morphisms}.
\end{align*}
\textit{The family \((M_i)\) satisfies axioms \SGARefLink{sga3:iv:para:4:6:3:item:a:01}{(a)}, \SGARefLink{sga3:iv:para:4:6:3:item:b:01}{(b)}, \SGARefLink{sga3:iv:para:4:6:3:item:c:01}{(c)},
\SGARefLink{sga3:iv:para:4:6:3:item:d:mathcal:t:01}{\((d_{T_i})\)}, and \SGARefLink{sga3:iv:para:4:6:3:item:e:mathcal:t:01}{\((e_{T_i})\)} of \SGARefLink{sga3:iv:para:4:6:3}{4.6.3}.}
\footnote{Editor's note: the original has been corrected by adding
surjectivity to \((M_3)\) and \((M_4)\); it is automatically satisfied in
the other two cases.}

Conditions \SGARefLink{sga3:iv:para:4:6:3:item:a:01}{(a)}, \SGARefLink{sga3:iv:para:4:6:3:item:b:01}{(b)}, and \SGARefLink{sga3:iv:para:4:6:3:item:c:01}{(c)} are classical (see EGA and \SGARefLink{sga3:iv:bib:sga:1}{SGA~1},
passim).\footnote{Editor's note: see EGA~I, 6.6.4 for "quasi-compact";
EGA~II, 6.1.5 for "finite"; EGA~IV\textsubscript{1}, 1.6.2 for "locally
of finite presentation"; EGA~IV\textsubscript{2}, 2.2.13 for "faithfully
flat"; and EGA~IV\textsubscript{4}, 17.3.3 for "\'{e}tale".}
By \SGARefLink{sga3:iv:proposition:6:3:1:item:iv:01}{6.3.1(iv)},\SGARefLink{sga3:iv:proposition:6:3:1:item:v:01}{(v)}, \((M_i)\) satisfies \SGARefLink{sga3:iv:para:4:6:3:item:d:mathcal:t:01}{\((d_{T_i})\)}. It remains to verify
\SGARefLink{sga3:iv:para:4:6:3:item:e:mathcal:t:01}{\((e_{T_i})\)}. It is enough to verify \SGARefLink{sga3:iv:para:4:6:3:item:e:mathcal:t:01}{\((e_{T_1})\)}, which implies all the
others; this follows from \SGARefLink{sga3:iv:bib:sga:1}{SGA~1}, VIII, \S\S4--5.

\medskip
\noindent\SGARefTarget{sga3:iv:corollary:6:3:3}\textbf{Corollary 6.3.3.}
\textit{Let \(X\) be a scheme and let \(R\) be an equivalence relation in
\(X\) of type \((M_i)\). Then \(R\) is \((M_i)\)-effective if and only if
the sheaf quotient of \(X\) by \(R\) for \(T_i\) is representable; in that
case it is represented by the quotient \(X/R\).}

This is \SGARefLink{sga3:iv:lemma:4:6:5}{4.6.5}.

\SGARefTarget{sga3:iv:subsection:6:4}\subsection*{6.4. Effectivity Conditions}
\addcontentsline{toc}{subsection}{\SGARefLink{sga3:iv:subsection:6:4}{6.4}. Effectivity Conditions}

We now seek families \((N)\) of morphisms satisfying axiom \SGARefLink{sga3:iv:subsection:4:7:item:f:mathcal:t:01}{\((f_T)\)} of
\SGARefLink{sga3:iv:subsection:4:7}{4.7}. Observe first that \SGARefLink{sga3:iv:subsection:4:7:item:f:mathcal:t:01}{\((f_{T_1})\)} implies \SGARefLink{sga3:iv:subsection:4:7:item:f:mathcal:t:01}{\((f_{T_i})\)}, so it is
enough to consider the \((\mathrm{fpqc})\) topology.

\medskip
\noindent\SGARefTarget{sga3:iv:lemma:6:4:1}\textbf{Lemma 6.4.1.}
\textit{The following families of morphisms satisfy axiom \SGARefLink{sga3:iv:subsection:4:7:item:f:mathcal:t:01}{\((f_{T_1})\)}
of \SGARefLink{sga3:iv:subsection:4:7}{4.7}, that is, they "descend by \((\mathrm{fpqc})\)":}
\begin{align*}
 (N)&:\ \text{open immersions};\\
 (N')&:\ \text{closed immersions};\\
 (N'')&:\ \text{quasi-compact immersions}.
\end{align*}

By \SGARefLink{sga3:iv:proposition:6:3:1:item:ii:01}{6.3.1(ii)}, it is enough to verify that these families descend for the
Zariski topology and along a faithfully flat quasi-compact morphism. The
first assertion is clear. For \((N)\), the second is \SGARefLink{sga3:iv:bib:sga:1}{SGA~1}, VIII, 4.4; for
\((N')\), it is loc. cit., 1.9; for \((N'')\), argue as in loc. cit., 5.5,
using the two preceding results.

\medskip
\noindent\SGARefTarget{sga3:iv:corollary:6:4:2}\textbf{Corollary 6.4.2.}
\textit{\SGARefLink{sga3:iv:lemma:6:4:1}{The same result} holds for quasi-compact open immersions.}

These results make the general results of \SGARefLink{sga3:iv:proposition:4:7:1}{4.7.1}, \SGARefLink{sga3:iv:proposition:4:7:2}{4.7.2}, \SGARefLink{sga3:iv:proposition:5:1:8}{5.1.8}, \SGARefLink{sga3:iv:proposition:5:3:1}{5.3.1}, and
so on applicable to the present situation. We state the first as an
example.

\medskip
\noindent\SGARefTarget{sga3:iv:corollary:6:4:3}\textbf{Corollary 6.4.3.}
\textit{\((=\SGARefLink{sga3:iv:proposition:4:7:1}{4.7.1}+\SGARefLink{sga3:iv:corollary:4:6:10}{4.6.10})\). Let \(X\) be a scheme and \(R\) an
equivalence relation in \(X\). Suppose that \(R\to X\) is faithfully flat
and quasi-compact, and that \(R\to X\times X\) is a closed immersion
(respectively, open, quasi-compact, or quasi-compact open). Then the sheaf
quotient \(X/R\) is the same for the \((\mathrm{fpqc})\) topology as for
the canonical topology. For every scheme \(S\),}
\[
(X/R)(S)=
\left\{
\begin{array}{l}
\text{closed (respectively, open, retrocompact, or open retrocompact)}
\\
\text{subschemes \(Z\) of \(X_S\), stable under \(R\times S\), such that}
\\
\text{\(Z\to S\) is faithfully flat and quasi-compact and}
\\
\text{\(R_Z\rightrightarrows Z\to S\) is exact}
\end{array}
\right\}.
\tag{6.4.3.1}\label{sga3:iv:corollary:6:4:3:eq:6:4:3:1:01}
\]
\footnote{Editor's note: recall that a subscheme \(Z\) of a scheme \(T\)
is called \emph{retrocompact} if the immersion \(Z\hookrightarrow T\) is
quasi-compact; see EGA~0\textsubscript{III}, 9.1.1.}

\SGARefTarget{sga3:iv:subsection:6:5}\subsection*{6.5. Principal Homogeneous Bundles}
\addcontentsline{toc}{subsection}{6.5. Principal Homogeneous Bundles}

We merely record the terminology:
\[
\begin{array}{c|l}
 \text{topology}&\text{principal homogeneous bundles}\\ \hline
 (\mathrm{fpqc})&\text{principal homogeneous bundles (without qualifier)}
 \\
 (\mathrm{\acute et})&\text{quasi-isotrivial principal homogeneous bundles}
 \\
 (\mathrm{\acute etf})&\text{locally isotrivial principal homogeneous bundles}
 \\
 (\mathrm{Zar})&\text{locally trivial principal homogeneous bundles}.
\end{array}
\tag{6.5.1}\label{sga3:iv:subsection:6:5:eq:6:5:1:01}
\]

\SGARefTarget{sga3:iv:subsection:6:6}\subsection*{6.6. Other Topologies}
\addcontentsline{toc}{subsection}{\SGARefLink{sga3:iv:subsection:6:6}{6.6}. Other Topologies}

Other topologies on the category of schemes are sometimes useful. One is
the \emph{global finite \'{e}tale topology}
\((\mathrm{\acute etfg})\), generated by the pretopology whose covering
families are the surjective families of finite \'{e}tale morphisms. It is
not finer than the Zariski topology. The corresponding principal
homogeneous bundles are called \emph{isotrivial}.

The position of this topology, and of the chaotic topology, in the lattice
is as follows:
\[
\begin{array}{ccccc}
  & & (\text{canonical}) & &\\
  & & \Big\downarrow & &\\[-0.25em]
  & & (\mathrm{fpqc}) & &\\
  & & \Big\downarrow & &\\[-0.25em]
  & & (\mathrm{fppf}) & &\\
  & & \Big\downarrow & &\\[-0.25em]
  & & (\mathrm{\acute et}) & &\\
  & & \Big\downarrow & &\\[-0.25em]
  & & (\mathrm{\acute etf}) & &\\
  & \swarrow & & \searrow &\\[-0.25em]
  (\mathrm{Zar}) & & & & (\mathrm{\acute etfg})\\
  & \searrow & & \swarrow &\\[-0.25em]
  & & (\text{chaotic}) & &
\end{array}
\tag{6.6.1}\label{sga3:iv:subsection:6:6:eq:6:6:1:01}
\]
\footnote{Editor's note: recall (cf. \SGARefLink{sga3:iv:scholium:4:4:2}{4.4.2}) that the chaotic topology is
the coarsest topology, defined by \(J(S)=\{S\}\) for every
\(S\in\operatorname{Ob}\mathscr C\).}

\SGARefTarget{sga3:iv:subsection:6:7}\subsection*{6.7. Homogeneous Spaces}
\addcontentsline{toc}{subsection}{\SGARefLink{sga3:iv:subsection:6:7}{6.7}. Homogeneous Spaces}

Let\footnote{Editor's note: the numbering that follows has been added.}
\(G\) be an \(S\)-group scheme, let \(X\) be an \(S\)-scheme equipped
with a left action of \(G\), and let
\[
  \Phi:G\times_SX\longrightarrow X\times_SX
\]
be the morphism of \(S\)-schemes defined on points by
\[
  (g,x)\longmapsto(gx,x).
\]
Recall (cf. \SGARefLink{sga3:iv:definition:5:1:0}{5.1.0} and III.\SGARefLink{sga3:iii:definition:0:1}{0.1}) that \(X\) is called a
\emph{formally principal homogeneous space} under \(G\) if the following
equivalent conditions hold:
\begin{enumerate}
  \item[(i)]\SGARefTarget{sga3:iv:subsection:6:7:item:i:01} for every \(T\to S\), the set \(X(T)\) is empty or is a
    principal homogeneous space under \(G(T)\);
  \item[(ii)]\SGARefTarget{sga3:iv:subsection:6:7:item:ii:01} \(\Phi\) is an isomorphism of \(S\)-functors;
  \item[(iii)]\SGARefTarget{sga3:iv:subsection:6:7:item:iii:01} \(\Phi\) is an isomorphism of \(S\)-schemes.
\end{enumerate}
The equivalence \SGARefLink{sga3:iv:subsection:6:7:item:i:01}{(i)}\(\Longleftrightarrow\)\SGARefLink{sga3:iv:subsection:6:7:item:ii:01}{(ii)} is clear, and
\SGARefLink{sga3:iv:subsection:6:7:item:ii:01}{(ii)}\(\Longleftrightarrow\)\SGARefLink{sga3:iv:subsection:6:7:item:iii:01}{(iii)} holds because
\(\mathscr C=(\mathbf{Sch}/S)\) is a full subcategory of
\(\widehat{\mathscr C}\).

The notion of a \emph{formally homogeneous space}, not necessarily
principal homogeneous, is obtained by requiring \(\Phi\) to be an
epimorphism in the category of sheaves for an appropriate topology
\(\mathcal T\). Indeed, requiring \(\Phi\) to be an epimorphism of
\(S\)-functors is equivalent to requiring that, for every \(T\to S\), the
set \(X(T)\) be empty or homogeneous (not necessarily principal
homogeneous) under \(G(T)\). This is too restrictive, as the following
simple example shows.

Let
\[
  S=\operatorname{Spec}R,\qquad G=\mathbf G_{m,R},\qquad
  X=\mathbf G_{m,R},
\]
and let \(G\) act on \(X\) by \(t\cdot x=t^2x\). Then \(\Phi\) is
surjective, finite, and \'{e}tale, and hence is an epimorphism in the
category of sheaves for the \((\mathrm{\acute etf})\) topology (and,
\textit{a fortiori}, an epimorphism of \(S\)-schemes). On the other hand,
the points \(1\) and \(-1\) of \(X(R)\) are not conjugate under any element
of \(G(R)\). Thus
\[
  G(R)\times X(R)\longrightarrow X(R)\times X(R)
\]
is not surjective.\footnote{Editor's note: this difficulty plainly comes
from the following fact. Let \(\mathscr C'\) be a full subcategory of
\(\widehat{\mathscr C}\) containing \(\mathscr C\), for example the
category \(\widetilde{\mathscr C}_{\mathcal T}\) of sheaves on
\(\mathscr C\) for a topology \(\mathcal T\) coarser than the canonical
topology. If \(f:X\to Y\) is a morphism in \(\mathscr C\), then the
implications
\[
 \text{\(f\) epimorphic in \(\widehat{\mathscr C}\)}
 \Longrightarrow
 \text{\(f\) epimorphic in \(\mathscr C'\)}
 \Longrightarrow
 \text{\(f\) epimorphic in \(\mathscr C\)}
\]
are generally strict.}
We are therefore led to \SGARefLink{sga3:iv:definition:6:7:1}{the following definition}.

\medskip
\noindent\SGARefTarget{sga3:iv:definition:6:7:1}\textbf{Definition 6.7.1.}
\textit{Let \(G\) be an \(S\)-group, let \(X\) be an \(S\)-scheme
equipped with an action of \(G\), and let \(\mathcal T\) be a topology on
\((\mathbf{Sch}/S)\) coarser than the canonical topology. The scheme \(X\)
is called a \emph{formally homogeneous space under \(G\)}, relative to
\(\mathcal T\), if the following equivalent conditions hold:}
\begin{enumerate}
  \item[(i)]\SGARefTarget{sga3:iv:definition:6:7:1:item:i:01} \textit{the morphism
    \(\Phi:G\times_SX\to X\times_SX\) is an epimorphism in the category of
    sheaves for \(\mathcal T\);}
  \item[(ii)]\SGARefTarget{sga3:iv:definition:6:7:1:item:ii:01} \textit{for every \(T\to S\) and every \(x,y\in X(T)\),
    there exist a \(\mathcal T\)-covering morphism \(T'\to T\) and an
    element \(g\in G(T')\) such that \(y_{T'}=g\cdot x_{T'}\).}
\end{enumerate}

\medskip
\noindent\SGARefTarget{sga3:iv:remark:6:7:2}\textbf{Remark 6.7.2.}
\textit{Condition \SGARefLink{sga3:iv:definition:6:7:1:item:i:01}{(i)} implies in particular that \(\Phi\) is a universal
effective epimorphism in \((\mathbf{Sch}/S)\), by \SGARefLink{sga3:iv:proposition:4:4:3}{4.4.3}. As is readily
seen, it follows that \(\Phi\) is surjective; cf. \SGARefLink{sga3:iv:definition:1:3}{1.3}, Editor's note (3).}

\medskip
\noindent\SGARefTarget{sga3:iv:propdef:6:7:3}\textbf{Proposition and Definition 6.7.3.}
\footnote{Editor's note: cf. \SGARefLink{sga3:iv:bib:ray70}{[Ray70]}, Definition VI.1.1.}
\textit{Let \(G\) be an \(S\)-group, let \(X\) be an \(S\)-scheme
equipped with an action of \(G\), and let \(\mathcal T\) be a topology on
\((\mathbf{Sch}/S)\) coarser than the canonical topology. The following
conditions are equivalent:}
\begin{enumerate}
  \item[(i)]\SGARefTarget{sga3:iv:propdef:6:7:3:item:i:01} \textit{\(X\) satisfies both of the following conditions:}
    \begin{enumerate}
      \item[(1)]\SGARefTarget{sga3:iv:propdef:6:7:3:item:1:01} \textit{the morphism
        \(\Phi:G\times_SX\to X\times_SX\) is covering, that is, \(X\) is a
        formally homogeneous \(G\)-space;}
      \item[(2)]\SGARefTarget{sga3:iv:propdef:6:7:3:item:2:01} \textit{the structural morphism \(X\to S\) is also
        covering, that is, locally for \(\mathcal T\) it has a section
        (cf. \SGARefLink{sga3:iv:remark:4:4:8:bis}{4.4.8 bis}, \SGARefLink{sga3:iv:remark:4:4:8:bis:item:b:01}{(b)}).}
    \end{enumerate}

  \item[(ii)]\SGARefTarget{sga3:iv:propdef:6:7:3:item:ii:01} \textit{Locally on \(S\) for the topology \(\mathcal T\),
    the scheme \(X\), with its \(G\)-action, is isomorphic to the
    quotient sheaf of \(G\) by a subgroup scheme \(H\), for
    \(\mathcal T\).
    In other words, there is a covering family \(\{S_i\to S\}\) such that
    every \(X\times_SS_i\) represents the sheaf quotient of
    \(G\times_SS_i\) by some subgroup scheme \(H_i\).}
\end{enumerate}
\textit{Under these conditions, \(X\) is called a
\emph{\(G\)-homogeneous space}, relative to \(\mathcal T\).}

\smallskip
\noindent\textit{Proof.}
Assume \SGARefLink{sga3:iv:propdef:6:7:3:item:ii:01}{(ii)}. Put
\[
  G_i=G\times_SS_i,\qquad X_i=X\times_SS_i.
\]
The scheme \(X_i\) has a section over \(S_i\), namely the composite of the
unit section \(\varepsilon_i:S_i\to G_i\) with the projection
\(\pi_i:G_i\to X_i=G_i/H_i\). It follows that \(X\to S\) is covering.

Moreover, \(\pi_i\) is covering, so \(\pi_i\times\pi_i\) is covering as
well, by \SGARefLink{sga3:iv:definition:4:2:3}{4.2.3} \SGARefLink{sga3:iv:definition:4:2:3:item:c:1:01}{(C~1)},\SGARefLink{sga3:iv:definition:4:2:3:item:c:2:01}{(C~2)}. There is a commutative
diagram\footnote{Editor's note: the Polo--Gille re-edition prints
\(G_i\times_{S_i}X_i\) at the upper right. We write
\(X_i\times_{S_i}X_i\), since \(\Phi_i\), as the base change of
\(\Phi:G\times_SX\to X\times_SX\), and \(\pi_i\times\pi_i\) both have
this codomain.}
\[
\begin{array}{ccccc}
  G_i\times_{S_i}X_i & \xrightarrow{\ \Phi_i\ } &
    X_i\times_{S_i}X_i\\
  \scriptstyle\operatorname{id}\times\pi_i\,\Big\uparrow
    & & \scriptstyle\pi_i\times\pi_i\,\Big\uparrow\\[-0.25em]
  G_i\times_{S_i}G_i & \xrightarrow[\sim]{\ \Psi_i\ } &
    G_i\times_{S_i}G_i ,
\end{array}
\tag{6.7.3.1}\label{sga3:iv:propdef:6:7:3:eq:6:7:3:1:01}
\]
where \(\Phi_i\) is obtained from \(\Phi\) by the base change \(S_i\to S\),
and \(\Psi_i\) is the isomorphism defined on points by
\[
  (g,g')\longmapsto(gg',g).
\]
The composite \((\pi_i\times\pi_i)\circ\Psi_i\) is covering; hence
\(\Phi_i\) is covering by \SGARefLink{sga3:iv:definition:4:2:3}{4.2.3} \SGARefLink{sga3:iv:definition:4:2:3:item:c:3:01}{(C~3)}. Thus \(\Phi\) is locally covering
and is therefore covering, by \SGARefLink{sga3:iv:definition:4:2:3}{4.2.3} \SGARefLink{sga3:iv:definition:4:2:3:item:c:5:01}{(C~5)}. This proves
\SGARefLink{sga3:iv:propdef:6:7:3:item:ii:01}{(ii)}\(\Longrightarrow\)\SGARefLink{sga3:iv:propdef:6:7:3:item:i:01}{(i)}.

Conversely, assume \SGARefLink{sga3:iv:propdef:6:7:3:item:i:01}{(i)}, and first suppose in addition that the structural
morphism \(X\to S\) has a section \(\sigma\). By EGA~I, 5.3.13,
\(\sigma\) is an immersion. Define \(H=G\times_XS\) by \SGARefLink{sga3:iv:propdef:6:7:3:eq:6:7:3:2:01}{the following
diagram}, whose two squares are cartesian:
\[
\begin{array}{ccccc}
  H & \longrightarrow & G
    & \xrightarrow{\ \operatorname{id}_G\boxtimes\sigma\ } & G\times_SX\\
  \Big\downarrow & & \scriptstyle\pi\,\Big\downarrow
    & & \scriptstyle\Phi\,\Big\downarrow\\[-0.25em]
  S & \xrightarrow{\ \sigma\ } & X
    & \xrightarrow{\ \operatorname{id}_X\boxtimes\sigma\ } & X\times_SX .
\end{array}
\tag{6.7.3.2}\label{sga3:iv:propdef:6:7:3:eq:6:7:3:2:01}
\]
Here \(\pi\), \(\operatorname{id}_G\boxtimes\sigma\), and
\(\operatorname{id}_X\boxtimes\sigma\) are the morphisms defined on points,
for \(T\to S\), \(g\in G(T)\), and \(x\in X(T)\), by
\[
 \pi(g)=g\cdot\sigma_T,\qquad
 (\operatorname{id}_G\boxtimes\sigma)(g)=(g,\sigma_T),\qquad
 (\operatorname{id}_X\boxtimes\sigma)(x)=(x,\sigma_T).
\]
Then \(\pi\) is covering, and \(H\) is a subgroup scheme of \(G\)
representing the stabilizer \(\operatorname{Stab}_G(\sigma)\) of \(\sigma\)
(cf. I, \SGARefLink{sga3:i:definition:2:3:3}{2.3.3}). Thus, for every \(T\to S\),
\[
  H(T)=\{g\in G(T)\mid g\cdot\sigma_T=\sigma_T\}.
\]

Let \(G/H\) denote the presheaf
\[
  T\longmapsto G(T)/H(T),
\]
and let \(a(G/H)\) be its associated sheaf for \(\mathcal T\). The
preceding construction gives a commutative diagram of presheaves equipped
with a \(G\)-action:
\[
\begin{array}{ccccc}
  G & \xrightarrow{\ \pi\ } & X\\
  \Big\downarrow & \nearrow_{\scriptstyle\bar\pi} &\\[-0.25em]
  G/H & &
\end{array}
\tag{6.7.3.3}\label{sga3:iv:propdef:6:7:3:eq:6:7:3:3:01}
\]
where \(\bar\pi\) is a monomorphism. Since \(\pi\) is covering,
\(\bar\pi\) is covering as well. By \SGARefLink{sga3:iv:corollary:4:3:12}{4.3.12}, \(\bar\pi\) therefore induces
an isomorphism
\[
  a(G/H)\xrightarrow{\sim}X.
\]
We have proved: if \(X\) is a \(G\)-homogeneous space such that \(X\to S\)
has a section \(\sigma\), then \(X\) represents the sheaf quotient \(G/H\),
where \(H=G\times_XS\) is the stabilizer of \(\sigma\).

In general, by hypothesis there is a covering family \(\{S_i\to S\}\) such
that every morphism
\[
  X_i=X\times_SS_i\longrightarrow S_i
\]
has a section \(\sigma_i\). Put \(G_i=G\times_SS_i\). The morphism
\[
  \Phi_i:G_i\times_{S_i}X_i\longrightarrow X_i\times_{S_i}X_i
\]
obtained from \(\Phi\) by base change along \(S_i\to S\) is still covering.

The preceding argument then gives
\[
  X_i\simeq G_i/H_i,
\]
where \(H_i\) is the stabilizer of \(\sigma_i\) in \(G_i\). This completes
the proof of \SGARefLink{sga3:iv:propdef:6:7:3:item:i:01}{(i)}\(\Longrightarrow\)\SGARefLink{sga3:iv:propdef:6:7:3:item:ii:01}{(ii)}.

\SGARefTarget{sga3:iv:section:bibliography}\section*{Bibliography}
\addcontentsline{toc}{section}{Bibliography}

\begin{description}
  \item[{[AS]}]\SGARefTarget{sga3:iv:bib:as} J.~Giraud, \emph{Analysis Situs}, S\'eminaire Bourbaki,
    Expos\'e 256, May 1963.

  \item[{[D]}]\SGARefTarget{sga3:iv:bib:d} J.~Giraud, \emph{M\'ethode de la descente}, M\'emoires de la
    Soci\'et\'e Math\'ematique de France, vol.~2 (1964), pp.~iii--viii and
    1--150.

  \item[{[MA]}]\SGARefTarget{sga3:iv:bib:ma} M.~Artin, \emph{Grothendieck Topologies}, mimeographed
    notes, Harvard, 1962.

  \item[{[SGA 1]}]\SGARefTarget{sga3:iv:bib:sga:1} \emph{S\'eminaire de G\'eom\'etrie Alg\'ebrique du
    Bois-Marie 1960--61: Rev\^etements \'etales et groupe fondamental},
    Lecture Notes in Mathematics 224 (1971); recomposed and annotated
    edition, Documents Math\'ematiques 3, Soci\'et\'e Math\'ematique de
    France, 2003.

  \item[{[SGA 4]}]\SGARefTarget{sga3:iv:bib:sga:4} \emph{S\'eminaire de G\'eom\'etrie Alg\'ebrique du
    Bois-Marie 1963--1964: Th\'eorie des topos et cohomologie \'etale des
    sch\'emas}, vols.~I--III, Lecture Notes in Mathematics 269, 270 (1972)
    and 305 (1973).

  \item[{[TDTE I]}]\SGARefTarget{sga3:iv:bib:tdte:i} A.~Grothendieck, \emph{Techniques de descente et
    th\'eor\`emes d'existence en g\'eom\'etrie alg\'ebrique I.
    G\'en\'eralit\'es. Descente par morphismes fid\`element plats},
    S\'eminaire Bourbaki, Expos\'e 190, December 1959.

  \item[{[Ray70]}]\SGARefTarget{sga3:iv:bib:ray70} M.~Raynaud, \emph{Faisceaux amples sur les sch\'emas en
    groupes et les espaces homog\`enes}, Lecture Notes in Mathematics 119,
    Springer-Verlag, 1970.\footnote{Editor's note: this reference has been
    added.}
\end{description}
